% limitations of our approach tf-idf features, linear models, etc. and potential future work to address these limitations (e.g., exploring other feature extraction methods, trying non-linear models, etc.)

% The remaining 8-9\% error rate likely reflects both feature limitations and genuine class ambiguity. Two complementary directions could address this: (a) richer feature representations such as word embeddings or contextual encodings (e.g., sentence transformers), which could better capture semantic similarity and polysemy that TF-IDF misses; and (b) non-linear models (kernel SVM, random forests, or neural classifiers) that could exploit more complex decision boundaries, as suggested by the overlapping clusters in the t-SNE plots. However, both directions carry trade-offs in interpretability, computational cost, and risk of overfitting on a dataset of this size, and would require careful tuning to confirm any gains over the linear baseline.


Our results point to two interconnected limitations. The first is the bag-of-words assumption underlying TF-IDF. Because the representation weights tokens purely by frequency with no awareness of context, the model cannot distinguish how a word is being used, only whether it appears. This directly explains the overlapping keywords errors we identified: an article containing terms strongly associated with both business and sci/tech will be pulled in both directions, and if the business-related terms happen to carry higher TF-IDF scores, the model will misclassify the article even if its subject matter is primarily technological. The same issue underlies the "simply difficult" category, where the balance of term frequencies failed to reflect the article's overall framing. Contextual encodings such as a fine-tuned transformer would address the two limitaitons; static word embeddings like Word2Vec would help with dealing with synonyms, but not with nuanced contextual meaning.

The second limitation is the linearity of our classifiers. The t-SNE projections show substantial overlap between the business and sci/tech clusters, which a linear decision boundary cannot cleanly resolve. Non-linear models could in principle learn more expressive boundaries in this region. However, our ablation results suggest that the representation is the primary bottleneck: the gains available from preprocessing changes were modest, which implies that the ceiling imposed by the bag-of-words assumption matters more than the choice of classifier. Switching to a transformer encoder would address both limitations simultaneously, at the cost of the interpretable feature analysis shown in Figure \ref{fig:linear_svm_top_features}.